% rn-animations-gestures.tex — Animated, LayoutAnimation, Reanimated, Gesture Handler;
% why UI-thread animation survives a busy JS thread.
% Sources (checked 2026-10-06):
%   reactnative.dev/docs/animations (native driver: non-layout properties only — transform,
%     opacity; not Flexbox/position) · /docs/layoutanimation (New Architecture: enabled
%     automatically, setLayoutAnimationEnabledExperimental is a no-op)
%   github.com/software-mansion/react-native-reanimated README + docs guides/migration-from-3.x
%     (4.x New Architecture only; worklet functions moved to react-native-worklets; runOnJS →
%     scheduleOnRN, runOnUI → scheduleOnUI, old names deprecated re-exports)
%   github.com/software-mansion/react-native-gesture-handler docs guides/upgrading-to-3,
%     fundamentals/reanimated-interactions + getting-started (hook API; renamed callbacks and
%     relations; UI-thread callbacks by default, runOnJS property); GitHub releases (3.0.0
%     2026-05-28, 3.3.0 latest 3.x; 2.33.0 on 2026-09-14)
% @labels: area=react-native kind=api platform=cross-platform topic=ui,performance discussion=197
% @tags: reanimated, worklets, usesharedvalue, useanimatedstyle, usenativedriver, animated-api, layoutanimation, gesture-handler, runonjs, withspring, ui-thread
\documentclass[8pt]{extarticle}
\usepackage{printup-sheet}

\lstdefinelanguage{TSSheet}{
  morekeywords={const,let,function,return,import,from,export,default,type,
    interface,if,else,new,true,false,null,undefined,async,await,
    useSharedValue,useAnimatedStyle,withSpring,withTiming,scheduleOnRN},
  sensitive=true, morecomment=[l]{//}, morecomment=[s]{/*}{*/},
  morestring=[b]", morestring=[b]', morestring=[b]`}

\tikzset{
  lane/.style={font=\scriptsize\bfseries, anchor=east, align=right},
  ok/.style={draw=sheetGreen!70!black, fill=sheetGreen!20, minimum height=3.2mm,
             minimum width=4.2mm, inner sep=0pt},
  miss/.style={draw=sheetRed, dashed, minimum height=3.2mm, minimum width=4.2mm, inner sep=0pt},
  busy/.style={draw=sheetRed, fill=sheetRed!18, minimum height=3.2mm, inner sep=0pt, font=\tiny},
  lbl/.style={font=\tiny, text=black!75, inner sep=1pt, align=center},
  sb/.style={box, font=\scriptsize, inner sep=1.5pt, minimum height=5mm},
}

\begin{document}

\sheettitle{RN animations \& gestures — Animated, Reanimated, RNGH}{react-native · folio}

\oneliner{Smooth animation means computing each frame \textbf{on the UI thread}, so a busy JS
thread cannot drop it. Animated's \texttt{useNativeDriver} does this for a pre-declared
animation of \emph{transform/opacity}; \textbf{Reanimated} runs your own \emph{worklets} on the
UI thread; \textbf{Gesture Handler} recognises touches natively and feeds them. iOS analogue:
Core Animation keeps running while your code blocks.}

\vspace{3pt}
\noindent\begin{tikzpicture}[sheet]
  % ── timeline: JS thread busy ──
  \node[font=\scriptsize\bfseries, anchor=west] at (-0.2,3.6) {JS thread busy for 100 ms — what the user sees (one box = one 16.7 ms frame)};
  \foreach \i in {0,...,11} \draw[sheetGrey!40] (1.3+\i*0.62,3.3) -- (1.3+\i*0.62,0.55);
  \node[lane] at (1.2,2.95) {JS thread};
  \node[ok, minimum width=10mm, anchor=west] at (1.32,2.95) {};
  \node[busy, minimum width=37mm, anchor=west] at (2.56,2.95) {JSON.parse + re-render 100 ms};
  \node[ok, minimum width=6mm, anchor=west] at (6.3,2.95) {};
  \node[lane] at (1.2,2.15) {JS-driven\\(Animated, no\\native driver)};
  \foreach \i in {0,1} \node[ok, anchor=west] at (1.34+\i*0.62,2.15) {};
  \foreach \i in {2,...,7} \node[miss, anchor=west] at (1.34+\i*0.62,2.15) {};
  \foreach \i in {8,...,11} \node[ok, anchor=west] at (1.34+\i*0.62,2.15) {};
  \node[lbl, text=sheetRed] at (4.4,1.75) {6 frames frozen, then a jump};
  \node[lane] at (1.2,1.1) {UI-thread\\(native driver /\\Reanimated)};
  \foreach \i in {0,...,11} \node[ok, anchor=west] at (1.34+\i*0.62,1.1) {};
  \node[lbl, text=sheetGreen!55!black] at (4.4,0.7) {every frame on time — but a tap \texttt{onPress} still waits for JS};
  % ── Reanimated + RNGH data flow ──
  \node[font=\scriptsize\bfseries, anchor=west] at (9.3,3.6) {Gesture $\to$ worklet $\to$ style, never touching JS};
  \draw[sheetBlue!60, dashed, rounded corners] (9.3,3.3) rectangle (16.5,1.65);
  \node[lbl, anchor=north west, text=sheetBlue] at (9.35,3.25) {UI thread (main)};
  \node[lbl, anchor=north east, text=sheetGreen!55!black] at (16.45,3.25) {view props set directly — no React render, no JS};
  \node[sb] (g) at (10.4,2.3) {native pan\\recogniser};
  \node[sb] (w) at (12.45,2.3) {\texttt{onUpdate}\\worklet};
  \node[sb, fill=sheetGreen!15, draw=sheetGreen!70!black] (sv) at (14.2,2.3) {shared\\value};
  \node[sb] (st) at (15.75,2.3) {animated\\style};
  \draw[flow] (g) -- (w); \draw[flow] (w) -- node[lbl, above]{\texttt{.value=}} (sv);
  \draw[flow] (sv) -- (st);
  \node[sb, draw=sheetOrange, fill=sheetOrange!10, minimum width=34mm] (js) at (13.3,0.75) {JS thread: React state, navigation};
  \draw[hot, dashed] (w.south) -- node[lbl, left]{\texttt{scheduleOnRN}} (w.south |- js.north);
  \draw[hot, dashed] (sv.south |- js.north) -- node[lbl, right]{\texttt{scheduleOnUI} · write \texttt{.value}} (sv.south);
\end{tikzpicture}

\begin{multicols}{2}
\raggedright

\section{Animated (core)}
\begin{itemize}
  \item \texttt{Animated.Value} (held in \texttt{useRef}) · \texttt{ValueXY}; drivers
        \texttt{timing} (duration + easing), \texttt{spring} (physical, interruptible),
        \texttt{decay} (velocity); compose \texttt{sequence}/\texttt{parallel}/\texttt{stagger}/\texttt{loop}.
  \item \texttt{interpolate(\{inputRange, outputRange,} \texttt{extrapolate:'clamp'\})} —
        one value drives many props (rotation, colour, opacity).
  \item \textbf{\texttt{useNativeDriver: true}} serialises the animation graph to native
        \emph{once}; frames then run on UI. Only \textbf{non-layout} props:
        \texttt{transform}, \texttt{opacity}. \textbf{Not} \texttt{width}, \texttt{height},
        \texttt{top}, \texttt{margin}, \texttt{flex} — they need Yoga layout every frame.
  \item \texttt{Animated.event(..., \{useNativeDriver:true\})} binds scroll offset
        natively — the classic collapsing header.
  \item \textbf{LayoutAnimation}: \texttt{configureNext(Presets.easeInEaseOut)} before
        a state change animates the \emph{next layout} — cheap, global, no gestures, not
        interruptible. Old-arch Android: \texttt{UIManager.setLayoutAnimation}\allowbreak\texttt{EnabledExperimental(true)}.
\end{itemize}

\section{Reanimated — how it wins}
\begin{itemize}
  \item A \textbf{worklet} (\texttt{'worklet'} directive, added by the Babel plugin
        for hook callbacks) is copied into a \textbf{second JS runtime on the UI
        thread}; captured variables are copied, not shared.
  \item \texttt{useSharedValue} — readable/writable from both runtimes via
        \texttt{.value}; writing it does \textbf{not} re-render React.
  \item \texttt{useAnimatedStyle(fn)} re-runs on UI when a shared value it reads
        changes and patches the native view — \textbf{any} prop, layout included.
  \item \texttt{withTiming}, \texttt{withSpring}, \texttt{withDecay}, \texttt{withSequence},
        \texttt{withRepeat}, \texttt{withDelay}; \texttt{useDerivedValue};
        \texttt{useAnimated}\allowbreak\texttt{ScrollHandler}; \texttt{cancelAnimation}.
  \item \texttt{scheduleOnRN(fn, ...args)} UI $\to$ JS; \texttt{scheduleOnUI(fn, ...args)} JS
        $\to$ UI (\texttt{react-native-worklets}; 3.x \texttt{runOnJS(fn)(args)} /
        \texttt{runOnUI} are deprecated re-exports in 4). Async hops — at the end of a
        gesture, never per frame.
  \item \textbf{Layout animations}: \texttt{entering=\{FadeIn\}}, \texttt{exiting=\{FadeOut\}},
        \texttt{layout=\{LinearTransition\}} on \texttt{Animated.View}.
  \item Reanimated 4 requires the New Architecture; worklets moved to
        \texttt{react-native-worklets}.
\end{itemize}

\columnbreak

\section{Example — swipe card, springs back}
\begin{lstlisting}[language=TSSheet]
const x = useSharedValue(0);
const pan = usePanGesture({                  // RNGH 3
  onUpdate: (e) => { x.value = e.translationX; }, // UI
  onDeactivate: (e) => {
    if (Math.abs(e.translationX) > 120) scheduleOnRN(dismiss);
    else x.value = withSpring(0, { velocity: e.velocityX });
  },
});
const style = useAnimatedStyle(() => ({
  transform: [{ translateX: x.value }],
  opacity: interpolate(Math.abs(x.value), [0, 200], [1, 0.3],
                       Extrapolation.CLAMP),
}));
return <GestureDetector gesture={pan}>
  <Animated.View style={[styles.card, style]} /></GestureDetector>;
\end{lstlisting}

\section{Gesture Handler 3 (hooks; 2.x still maintained)}
\begin{itemize}
  \item \textbf{Native recognisers} (UIGestureRecognizer on iOS), replacing the JS responder
        system — they work while JS is busy. Root: \texttt{GestureHandlerRootView}.
  \item A config object per hook: \texttt{usePanGesture(\{…\})}, \texttt{useTapGesture},
        \texttt{useLongPressGesture}, \texttt{usePinchGesture}, \texttt{useRotationGesture},
        \texttt{useFlingGesture} $\to$ \texttt{<GestureDetector gesture=\{g\}>}. Callbacks
        \texttt{onBegin} · \texttt{onActivate} · \texttt{onUpdate} · \texttt{onDeactivate}
        (v2 \texttt{Gesture.Pan().onStart/.onEnd}; \texttt{onChange} folded into \texttt{onUpdate}).
  \item \texttt{useSimultaneousGestures(pinch, rotate)} — both active ·
        \texttt{useExclusiveGestures(doubleTap, tap)} — first that succeeds, in order ·
        \texttt{useCompetingGestures} (v2 \texttt{Race}). Across components:
        \texttt{requireToFail} (UIKit \texttt{require(toFail:)}), \texttt{simultaneousWith},
        \texttt{block}.
  \item Hook-API callbacks run on the \textbf{UI thread} by default (Reanimated
        integration); \texttt{runOnJS: true} moves them to JS.
\end{itemize}

\section{Traps}
\begin{itemize}
  \trap{``Native driver makes everything smooth'' — not layout props, and a JS
        \texttt{onPress}/\texttt{setState} is still late.}
  \trap{Calling a plain JS function from a worklet without \texttt{scheduleOnRN} crashes;
        reading \texttt{sv.value} in render gives a stale value.
        \texttt{scheduleOnRN(setState, v)} per frame rebuilds the JS bottleneck.}
  \trap{Animating \texttt{width}/\texttt{height} costs a layout pass per frame even on UI —
        prefer \texttt{transform: scale} where the look allows.}
\end{itemize}

\section{Remember}
\textbf{Frames on UI, decisions on JS}: gesture $\to$ worklet $\to$ shared value $\to$
style; cross to JS only at the end.


\end{multicols}

\noindent{\footnotesize\color{sheetGrey}\textit{Related:} rn-performance · Core Animation
(render server) · UIKit touches (UIGestureRecognizer) · rn-architecture (JSI, Fabric)}

\end{document}
